\section{Zoo Consumption and PQ-Rooted Omnichain Settlement}
\label{sec:l3-settlement}

The execution proof of Sections~\ref{sec:freivalds}--\ref{sec:families} makes an admitted
AI computation accountable. It does not create a new mining authority wherever the proof
is verified. The canonical topology has one authority and three downstream consumers:

\begin{enumerate}[leftmargin=1.4em]
\item \textbf{A-Chain PoAI consensus} admits the demand-backed job, fixes its execution
and proof policy, derives the post-commitment challenge, resolves the evidence to final or
rejected, consumes the global work nullifier, and authorizes one reward receipt.
\item \textbf{Z-Chain P3Q settlement} batches finalized A-Chain receipt transitions into a
Plonky3-derived, pairing-free STARK/FRI proof. It proves compression and settlement of an
A-Chain decision; it does not admit raw work or calculate a separate subsidy.
\item \textbf{Lux Quasar} orders and finalizes the A-Chain and Z-Chain state under the
actually activated post-quantum validator policy.
\item \textbf{Zoo and other destination applications} verify the PQ-finalized, Z-settled
A-Chain receipt, enforce destination and recipient binding, record local consumption, and
release the authorized representation exactly once.
\end{enumerate}

This separation resolves an otherwise fatal supply ambiguity. ``Mine into Zoo'' means an
A-Chain job names a Zoo destination before execution. It does not mean Zoo accepts a raw
transcript, runs another PoAI vote, or owns another spent set. One job consumes one global
nullifier on A-Chain and can yield one receipt, regardless of where the receipt is paid.

\subsection{What Zoo contributes}

Zoo remains economically and scientifically important without duplicating consensus. Zoo
applications can fund inference, embedding, evaluation, and supported training jobs; curate
models and datasets; define public-good reward programs; and use finalized execution
receipts as reproducibility artifacts. A scientific result can cite the committed model,
input, kernel profile, graph, output, and transcript root rather than asking readers to
trust that an opaque provider ran the advertised computation.

The same mechanism covers the heavy arithmetic of several task families:

\begin{itemize}[leftmargin=1.4em]
\item a proven inference binds the declared model and input to its graph and output;
\item a proven embedding binds a representation to its declared source computation;
\item a supported training step is a longer deterministic graph whose forward, backward,
and update matmuls use the same local verifier; and
\item a governed evaluation job may combine objective execution evidence with separately
attributable human or model judgment, while judgment alone cannot authorize mining subsidy.
\end{itemize}

Cryptography proves execution of the specified job, not universal truth, dataset rights,
or scientific importance. Demand-backed admission supplies the protocol meaning of useful;
governance and peer review remain responsible for the quality of the questions and data.

\subsection{Receipt path and exactly-once payout}

The exported object is a finalized receipt, never a raw \texttt{ComputeProof}, TEE quote,
miner signature, or Freivalds opening. The receipt binds the A-Chain network and policy
epoch, job, accepted execution commitment, amount, destination chain, recipient, and
nullifier. Z-Chain settlement binds that receipt transition to an A-Chain receipt root;
Quasar finality authenticates the A/Z roots and heights; a destination verifies inclusion
and records consumption of that exact receipt.

A destination-local consumed set prevents replay of an already valid authorization. It
does not judge whether work occurred. That judgment exists only on A-Chain, which is what
makes a capped issuance schedule compatible with omnichain payout.

\subsection{P3Q activation boundary}

Before the PoAI AIR and verifier are activated, Z-Chain may carry only an authenticated,
already-final A-Chain receipt root. It must not present that checkpoint as a succinct proof
of the underlying work. P3Q authority activates only after the AIR, parameters, recursion
policy, A-to-Z state-transition binding, verifier, resource bounds, node wiring, and audits
are complete. A shadow phase may compare candidate proofs against already-final receipts,
but those proofs have no mint authority. Staging is deployment safety, not a weaker proof
tier marketed to users.

The result is a clean omnichain contract: useful work is requested anywhere, executed on
any conforming hardware, judged once by A-Chain PoAI, compressed by Z-Chain P3Q, finalized
by Lux Quasar, and paid exactly once at the destination selected before mining.
